% ============================================================================
\clearpage
\aportacion[Escenarios]{Jacqueline Sarmiento Cervantes}{Matrícula A01795863 · Escenarios 1 y 2 · Personas: Laura Méndez y Diego Hernández}

% ---------------------------------------------------------------- Escenario 1
\section{Escenario 1 · Validar una cifra de cartera antes de una reunión}\label{sec:escenarios}

\escenariohead{imagenes/persona_01_laura.png}
  {Validar una cifra de cartera antes de una reunión}
  {Laura Méndez · Perfil operativo, consulta rápida · Persona primaria del producto}
  {Rol en el sistema: operativo · Equipo de escritorio · Urgencia alta, menos de treinta minutos}

\elemento{Definición del usuario}{audiencia objetivo, necesidades y objetivos, pain points}

Laura Méndez, 34 años, especialista de Operaciones Financieras con ocho años en el sector y
tres en su función actual. Nivel técnico intermedio: domina Excel y consulta sistemas
internos, pero no escribe consultas en lenguaje SQL ni programa. Pertenece al perfil
operativo de consulta rápida, el más numeroso de la audiencia del portal y el que la
actividad anterior designó como persona primaria del producto, por ser el denominador común
de todos los demás perfiles.

\textbf{\textcolor{uxnavy}{Necesidades y objetivos.}} Encontrar una cifra puntual y confirmar
su vigencia en pocos minutos. Responder las solicitudes que recibe con una fuente
identificable y una definición clara detrás. Resolver por sí misma las tareas frecuentes, sin
depender de que otra área conteste. Conservar las consultas habituales para no reconstruirlas
cada vez que se las piden.

\textbf{\textcolor{uxnavy}{Pain points relevantes para esta tarea.}} Recuerda nombres de
carpetas y sistemas, pero los accesos y las ubicaciones cambian sin que ella se entere.
Encuentra variables de nombre parecido y no sabe cuál corresponde a la solicitud que le
hicieron. Termina escribiendo a un colega para confirmar el responsable, la fecha o la fuente,
y el ritmo de su trabajo pasa a depender de la disponibilidad de esa persona. Los filtros
técnicos y los campos avanzados de las herramientas actuales estorban en una consulta que
debería ser simple.

\elemento{Descripción de la tarea}{qué quiere conseguir y pasos para lograrlo}

Son las 11:10. La subdirección de Laura tiene reunión a las 12:00 y le pide confirmar el saldo
de cartera vigente de un producto de crédito al cierre del mes anterior, acompañado de la
definición oficial de la variable y de la referencia de la fuente. La solicitud llega con esa
exigencia adicional porque en la reunión anterior dos áreas presentaron cifras distintas para
el mismo concepto y nadie pudo explicar de dónde salía cada una.

Laura no sabe de memoria qué área es propietaria de esa variable ni si la clasificación
cambió durante el último trimestre. Su meta no es únicamente obtener un número: es poder
responder de dónde salió si alguien pregunta en la reunión. Hoy una tarea como esta le toma
entre treinta y cuarenta y cinco minutos, y la mayor parte de ese tiempo transcurre esperando
la respuesta de un colega \emph{(estimación del equipo a partir de su experiencia
profesional; magnitud pendiente de medición con el instrumento diseñado en la actividad
anterior)}.

\figuraescena{imagenes/escena_laura.png}
  {Contexto del escenario 1: la validación ocurre en el escritorio, con la reunión
   aproximándose. Ilustración generada con inteligencia artificial.}{fig:escena-laura}

\textbf{\textcolor{uxnavy}{Pasos para lograrlo.}} El siguiente método describe la tarea sin
suponer ningún diseño concreto:

\begin{uxpreg}
\item Identificar el concepto exacto que le solicitaron, con sus posibles sinónimos.
\item Localizar qué fuente institucional contiene ese concepto.
\item Confirmar que la fuente es la vigente y no una versión retirada.
\item Leer la definición oficial de la variable y sus reglas de uso.
\item Identificar al área propietaria por si necesita respaldo posterior.
\item Acotar el dato al periodo solicitado.
\item Verificar que la cifra es consistente con lo que se reportó anteriormente.
\item Obtener el valor junto con su referencia de origen y su fecha de actualización.
\item Entregar la respuesta en el formato que su subdirección espera.
\item Dejar la consulta guardada para la próxima vez que se la pidan.
\end{uxpreg}

\elemento{Identificación de las interacciones clave}{puntos de contacto entre Laura y el portal}

Laura entra al portal con su cuenta institucional y el sistema le presenta su
\textbf{espacio de trabajo operativo}, con el buscador al centro y sus consultas guardadas a
un costado. Escribe el concepto en lenguaje de negocio, sin conocer el nombre técnico del
campo.

El \textbf{buscador unificado} devuelve tres resultados ordenados. Cada tarjeta muestra, antes
de que ella tenga que preguntarlo, la fuente, el área propietaria, la periodicidad y la fecha
de última actualización. Una de las tarjetas lleva una marca de versión retirada, con enlace
a la vigente: ese aviso resuelve por sí solo el riesgo que originó la solicitud.

Laura abre la ficha del \textbf{catálogo de datos} y lee la definición oficial, la nota de
conocimiento tribal que dejó el área propietaria —\emph{a partir de marzo la clasificación
separa el saldo vencido}— y el nombre de Roberto Valdez como responsable. No necesita
escribirle.

Aplica el \textbf{filtro de periodo} sobre el resultado; la vista conserva el contexto de la
búsqueda y no la obliga a reconstruirla. Con una duda sobre el alcance de la variable, escribe
la pregunta al \textbf{asistente conversacional} desde el mismo panel. El asistente muestra en
pantalla la tarjeta de la herramienta que está consultando —\emph{consultando el catálogo de
créditos}— antes de responder, y cada cifra de la respuesta cita la fuente de la que
proviene.

Finalmente \textbf{copia el dato con su referencia de origen y su fecha}, en un formato que
puede pegar directamente en el correo, y \textbf{guarda la consulta} en su espacio con un
nombre propio. Son las 11:23.

\elemento{Descripción del resultado deseado}{resultado de la tarea y sensación al completarla}

Laura responde la solicitud treinta y siete minutos antes de la reunión, con la cifra, la
definición oficial, el propietario y la fecha de corte. Si en la reunión alguien presenta otro
número, ella puede explicar de dónde salió el suyo y cuándo se actualizó.

La sensación que busca no es alivio por haber terminado a tiempo, sino \textbf{respaldo}: la
confianza de haber respondido con algo que resiste una pregunta incómoda delante de su
subdirección. Es una diferencia importante para el diseño, porque un producto que solo
optimiza la velocidad no produce esa sensación; la produce el hecho de que el resultado venga
acompañado de su procedencia. La siguiente vez que le pidan lo mismo, la consulta ya está
guardada y el trabajo se reduce a abrirla y actualizar el periodo.

\elemento{Excepciones}{situaciones poco frecuentes con consecuencias de diseño}

\begin{uxlist}
\lead{La fuente no responde}{el silo de créditos está en mantenimiento. El portal muestra el
  último valor disponible con su fecha y una marca visible de dato no vigente, en lugar de una
  pantalla en blanco o un error técnico. Laura decide con información suficiente si espera o
  responde con la salvedad.}
\lead{No hay una única respuesta}{la búsqueda devuelve dos variables legítimas con
  definiciones distintas. En vez de elegir por ella, el portal las muestra lado a lado con sus
  diferencias y sus propietarios. Laura escala con evidencia en lugar de adivinar.}
\lead{La definición cambió esta semana}{la ficha del catálogo muestra el cambio, su fecha y la
  versión anterior, y Laura menciona el cambio en su respuesta. Este caso es el punto de
  encuentro con el escenario 4, donde el mismo evento se observa desde el lado de quien lo
  publica.}
\end{uxlist}

% ---------------------------------------------------------------- Escenario 2
\clearpage
\section{Escenario 2 · Cruzar tres fuentes y dejar el análisis reproducible}

\escenariohead{imagenes/persona_02_diego.png}
  {Cruzar tres fuentes y dejar el análisis reproducible}
  {Diego Hernández · Perfil analista de datos, profundidad}
  {Rol en el sistema: analista · Escritorio y modalidad híbrida · Urgencia media, plazo de tres días}

\elemento{Definición del usuario}{audiencia objetivo, necesidades y objetivos, pain points}

Diego Hernández, 29 años, analista de datos financieros con cinco años de experiencia en
análisis, automatización y visualización. Nivel técnico avanzado: trabaja con SQL, Python,
Excel y herramientas de inteligencia de negocio. Pertenece al perfil de analista de datos, que
necesita combinar fuentes, aplicar filtros complejos, revisar metadatos y extraer datos crudos
para procesarlos en sus propias herramientas.

\textbf{\textcolor{uxnavy}{Necesidades y objetivos.}} Descubrir fuentes relacionadas y
cruzarlas con criterios consistentes. Guardar consultas, filtros y versiones para poder
repetir el análisis meses después. Obtener datos crudos en el formato que corresponda al
volumen: archivo cuando son miles de filas, conexión programática cuando son millones.
Identificar cobertura, periodicidad, propietario y limitaciones de cada fuente antes de
modelar sobre ella.

\textbf{\textcolor{uxnavy}{Pain points relevantes para esta tarea.}} Las fuentes usan
identificadores, catálogos y granularidades distintas, de modo que cruzarlas exige un trabajo
de conciliación que no aporta valor analítico. Las definiciones viven dispersas en correos,
archivos y en la memoria de especialistas. Algunos flujos de descarga no soportan archivos
pesados o exigen pasos manuales. Y reconstruir qué versión, qué filtro y qué transformación
produjeron un resultado es difícil, lo que le impide defender su propio análisis semanas
después.

\elemento{Descripción de la tarea}{qué quiere conseguir y pasos para lograrlo}

El área de riesgos le pide estimar la exposición conjunta por contraparte cruzando la cartera
de crédito, las posiciones de liquidez y las operaciones de derivados, con corte al último día
del trimestre. El entregable no es solo el número: el área quiere poder repetir el ejercicio
el trimestre siguiente sin volver a pedirle a Diego que reconstruya el procedimiento.

El reto real no es el cálculo, sino la conciliación. Diego sabe por experiencia que las tres
fuentes identifican a la contraparte de forma distinta, que la periodicidad de liquidez es
diaria mientras que la de crédito es mensual, y que un cruce ingenuo produciría un número
plausible y equivocado. Su meta es entregar una cifra que pueda explicar y un procedimiento
que otra persona pueda ejecutar.

\figuraescena{imagenes/escena_diego.png}
  {Contexto del escenario 2: el trabajo analítico transcurre entre varias fuentes abiertas a
   la vez y exige conciliarlas antes de cruzarlas. Ilustración generada con inteligencia
   artificial.}{fig:escena-diego}

\textbf{\textcolor{uxnavy}{Pasos para lograrlo.}}

\begin{uxpreg}
\item Localizar las tres fuentes que contienen la información requerida.
\item Revisar la definición, la periodicidad y el propietario de cada una.
\item Verificar qué identificador permite relacionarlas y con qué nivel de confianza.
\item Comprobar si la granularidad y las fechas de corte son compatibles.
\item Probar el cruce sobre una muestra reducida antes de procesar el volumen completo.
\item Revisar los valores atípicos que aparezcan y decidir si son error o realidad.
\item Ejecutar la extracción completa sin bloquear su sesión de trabajo.
\item Documentar filtros, fechas y transformaciones aplicadas.
\item Guardar la consulta de forma que pueda volver a ejecutarse con otro corte.
\item Compartir con el área solicitante la referencia del procedimiento, no solo el resultado.
\end{uxpreg}

\elemento{Identificación de las interacciones clave}{puntos de contacto entre Diego y el portal}

Diego parte del \textbf{catálogo de datos} y busca por tema en lugar de por nombre de tabla.
La ficha de la cartera de crédito le muestra las \textbf{fuentes relacionadas}, y ahí descubre
el conjunto de derivados, que no conocía porque pertenece a otra dirección. Ese
descubrimiento —que hoy dependería de que alguien se lo mencionara— es en sí mismo un
resultado del portal.

Compara las tres fichas y confirma en los metadatos la periodicidad, la fecha de corte y el
campo que actúa como identificador común de contraparte. El \textbf{explorador analítico} le
permite construir el cruce con filtros sobre las tres fuentes y le advierte, antes de
ejecutar, que la granularidad de liquidez es diaria y requiere una regla de agregación
explícita. Diego elige el cierre de mes y la decisión queda registrada como parte de la
consulta.

Ejecuta primero sobre una \textbf{muestra} de una contraparte conocida y verifica el resultado
contra un cálculo propio. Cuando el volumen completo supera el umbral de proceso inmediato, el
portal ofrece \textbf{exportar en segundo plano}: Diego lanza el trabajo, sigue navegando y
recibe el aviso con el enlace de descarga cuando termina, sin mantener una pestaña abierta ni
vigilar una barra de progreso.

Antes de cerrar, \textbf{guarda la consulta con nombre y versión} y comparte su referencia con
el área de riesgos. Cualquiera con permiso sobre las mismas fuentes puede abrir esa referencia
y ver exactamente qué filtros, qué cortes y qué regla de agregación produjeron el número.

\elemento{Descripción del resultado deseado}{resultado de la tarea y sensación al completarla}

Diego entrega la exposición conjunta por contraparte y, junto con ella, un procedimiento
reejecutable. El área de riesgos puede repetir el ejercicio el trimestre siguiente cambiando
únicamente la fecha de corte.

La sensación que busca es \textbf{control}: no la satisfacción de haber terminado, sino la
tranquilidad de poder explicarle a otra persona de dónde salió cada cifra y de saber que su
análisis no depende de que él recuerde lo que hizo. Para un perfil como el suyo, un producto
que entrega el dato pero no la trazabilidad del dato resuelve la mitad del problema y crea la
otra mitad.

\elemento{Excepciones}{situaciones poco frecuentes con consecuencias de diseño}

\begin{uxlist}
\lead{El volumen excede lo razonable para un archivo}{el portal propone la vía programática en
  lugar de generar una descarga de varios gigabytes, y ofrece la consulta ya formulada para
  ser consumida desde sus propias herramientas.}
\lead{Una de las tres fuentes está en mantenimiento}{el explorador no falla en silencio:
  informa qué fuente no está disponible, muestra el resultado parcial marcado como incompleto
  y permite reanudar el cruce cuando el servicio vuelva.}
\lead{Los valores atípicos son reales}{el portal permite anotar el hallazgo sobre la fuente
  para que el propietario del dato lo revise, de modo que la observación de Diego no se pierda
  en un correo y beneficie a quien consulte la misma fuente después.}
\end{uxlist}

% ------------------------------------------------------ Journey individual 1
\clearpage
\section{Journey map individual · Diego Hernández}

\textit{Elaborado por Jacqueline Sarmiento Cervantes sobre el escenario 2.}

El primero de los cuatro mapas de este documento despliega en el tiempo el objetivo del
escenario 2: \textbf{cruzar créditos, liquidez y derivados por contraparte y dejar el análisis
reejecutable}. Cubre desde que llega el encargo hasta que el procedimiento queda guardado y
compartido, es decir, también el antes y el después del contacto con el producto.

\textbf{Qué revela el mapa que el escenario no mostraba.} El escenario describe una tarea
razonablemente fluida. El mapa muestra que la curva emocional tiene dos caídas y que ninguna
de las dos está donde uno esperaría: no en la extracción de datos, que es la parte
técnicamente pesada, sino en la etapa 1, cuando Diego no puede estimar si el cruce es viable
sin abrir cada fuente, y en la etapa 3, cuando descubre que la granularidad de liquidez y la
de crédito no coinciden. Ambas son etapas de \emph{decisión}, no de ejecución, y hoy se
resuelven con conocimiento acumulado que el portal puede sustituir con metadatos.

\textbf{Momento de la verdad.} La etapa 2, el descubrimiento de una fuente relacionada que
Diego no conocía. Es el momento en que el portal deja de ser un lugar donde bajar archivos y
pasa a ser un lugar donde se encuentra algo que no se estaba buscando. Si esa etapa no
funciona, el resto del recorrido sigue siendo útil pero deja de ser diferencial.

\textbf{Consecuencia para el diseño.} Dos de las seis oportunidades del mapa —cobertura y
periodicidad visibles antes de empezar, y aviso de granularidad incompatible antes de
ejecutar— son advertencias tempranas, no funciones nuevas. Cuestan poco y actúan sobre las dos
caídas de la curva.

\figurapanoramica{figuras/journey_diego.pdf}
  {Journey map individual de Diego Hernández. Objetivo: cruzar créditos, liquidez y derivados
   por contraparte y dejar el análisis reejecutable. Elaboración propia.}{fig:journey-diego}

% ============================================================================
\clearpage
\aportacion[Escenarios]{Alexandro Mayoral Terán}{Matrícula A01795899 · Escenarios 3 y 4 · Personas: Arturo Castañeda y Roberto Valdez}

% ---------------------------------------------------------------- Escenario 3
\section{Escenario 3 · Entender una señal de riesgo antes de exponer ante la Junta}

\escenariohead{imagenes/persona_06_arturo.png}
  {Entender una señal de riesgo antes de exponer ante la Junta}
  {Arturo Castañeda · Perfil directivo, abstracción}
  {Rol en el sistema: directivo · Tableta institucional · Urgencia crítica, quince minutos}

\elemento{Definición del usuario}{audiencia objetivo, necesidades y objetivos, pain points}

Arturo Castañeda, 55 años, Director General de Estabilidad Financiera, economista con
veinticinco años de trayectoria en el sector público y financiero. Nivel técnico básico en
herramientas: interpreta gráficas y tableros con soltura, pero no escribe consultas ni
manipula bases de datos. Pertenece al perfil directivo, que necesita indicadores consolidados,
tendencias y señales de riesgo para supervisar resultados sin recorrer cada base por separado.

\textbf{\textcolor{uxnavy}{Necesidades y objetivos.}} Ver el estado general de la estabilidad
financiera de un vistazo. Identificar de inmediato anomalías, picos o tendencias
preocupantes. Tener la certeza de que el número que aparece en la pantalla es la versión
final, oficial y validada por sus gerentes.

\textbf{\textcolor{uxnavy}{Pain points relevantes para esta tarea.}} Le incomoda que en una
reunión un director presente un número y otro presente uno distinto porque usaron cortes de
fecha diferentes. Cuando un indicador se ve alarmante, tiene que hacer varias llamadas a sus
gerentes para entender la causa del movimiento, y el tiempo es su recurso más escaso. Recibe
reportes con demasiadas tablas y datos crudos que no le ayudan a entender el panorama.

\elemento{Descripción de la tarea}{qué quiere conseguir y pasos para lograrlo}

Faltan quince minutos para que Arturo entre a presentar ante la Junta de Gobierno. Va a exponer
la situación de riesgo del sistema bancario del mes y dispone de cinco minutos de intervención.
Quiere confirmar que los cuatro indicadores clave que va a presentar siguen siendo los mismos
que revisó ayer y, si alguno se movió, entender por qué antes de que se lo pregunten.

Su meta no es analizar: es \emph{poder responder}. La diferencia es sustantiva para el diseño.
Arturo no necesita acceso a los datos crudos; necesita que cada cifra que muestre venga
acompañada de la explicación que él daría si un subgobernador levantara la mano.

\figuraescena{imagenes/escena_arturo.png}
  {Contexto del escenario 3: la consulta ocurre de pie, en tableta y minutos antes de una
   reunión de alto nivel. Ilustración generada con inteligencia artificial.}{fig:escena-arturo}

\textbf{\textcolor{uxnavy}{Pasos para lograrlo.}}

\begin{uxpreg}
\item Revisar el estado general de los indicadores clave del periodo.
\item Detectar cuáles se movieron respecto de la última revisión.
\item Determinar si el movimiento es relevante o es ruido esperable.
\item Obtener la explicación del movimiento sin recurrir a una llamada.
\item Confirmar que la cifra es la versión oficial y está conciliada.
\item Verificar la fecha de corte que se está mostrando.
\item Quedarse con una frase clara para explicar el movimiento si se lo preguntan.
\end{uxpreg}

\elemento{Identificación de las interacciones clave}{puntos de contacto entre Arturo y el portal}

Arturo abre el portal desde su \textbf{tableta institucional} y accede a su \textbf{espacio
directivo}, que muestra los cuatro indicadores clave de riesgo bancario del mes con su
tendencia. La vista está preparada para lectura, no para exploración: pocas cifras, grandes,
con su variación respecto del periodo anterior.

Nota un alza en uno de los indicadores. Toca la gráfica y el portal despliega una
\textbf{tarjeta de contexto} escrita por el propietario del dato: \emph{el alza corresponde al
ajuste estacional aprobado la semana pasada}. Esa tarjeta es el sustituto exacto de las tres
llamadas que hoy tendría que hacer.

La misma vista muestra el \textbf{sello de conciliación} del indicador —quién lo validó y
cuándo— y la \textbf{fecha de corte} de forma visible, de manera que Arturo sabe con certeza
qué versión está presentando. Un nivel más de \textbf{profundización} le permite ver la
composición del indicador por tipo de institución, sin salir del tablero ni entrar a una
herramienta distinta.

Antes de guardar la tableta pide al \textbf{asistente} un resumen de la vista actual. El
asistente redacta tres líneas con el movimiento, su causa declarada y la fecha de corte,
citando la fuente de cada cifra. Arturo bloquea la tableta y entra a la reunión.

\elemento{Descripción del resultado deseado}{resultado de la tarea y sensación al completarla}

Arturo entra a la Junta sabiendo qué va a decir sobre el único indicador que se movió, con la
causa del movimiento y la fecha de corte presentes. Si alguien pregunta de dónde salió la
cifra, tiene la respuesta.

La sensación buscada es \textbf{certeza}, y es de un tipo particular: no la seguridad de estar
en lo correcto, sino la de no quedar expuesto. Para un perfil directivo el costo de un número
equivocado no es analítico sino institucional, y el producto debe reconocerlo. Un tablero
rápido pero sin sello de conciliación ni fecha de corte visible resolvería el problema
equivocado.

\elemento{Excepciones}{situaciones poco frecuentes con consecuencias de diseño}

\begin{uxlist}
\lead{El indicador todavía no está conciliado}{el portal lo declara con una marca explícita en
  lugar de mostrar un valor provisional como si fuera definitivo. Arturo prefiere una cifra
  marcada como preliminar a una cifra que después tenga que corregir en público.}
\lead{No hay tarjeta de contexto para el movimiento}{el portal ofrece solicitar la explicación
  al propietario del dato con un toque, y la solicitud queda registrada con hora, de modo que
  la ausencia de contexto se vuelve visible y accionable en lugar de silenciosa.}
\lead{La conexión es intermitente durante un traslado}{la vista directiva conserva la última
  lectura con su fecha y hora de descarga, claramente identificada como tal, para que Arturo
  nunca presente un dato creyendo que es de hoy cuando es de ayer.}
\end{uxlist}

% ---------------------------------------------------------------- Escenario 4
\clearpage
\section{Escenario 4 · Publicar una nueva versión del catálogo sin romper el trabajo de nadie}

\escenariohead{imagenes/persona_03_roberto.png}
  {Publicar una nueva versión del catálogo sin romper el trabajo de nadie}
  {Roberto Valdez · Perfil propietario de datos, gobierno y calidad}
  {Rol en el sistema: analista con atribuciones de gobierno · Escritorio · Urgencia baja, ventana de dos semanas}

\elemento{Definición del usuario}{audiencia objetivo, necesidades y objetivos, pain points}

Roberto Valdez, 52 años, subgerente de Información de Estados Financieros y Crédito, contador
con dieciocho años en la institución y experto en catálogos y metodologías crediticias. Nivel
técnico intermedio: bases de datos relacionales, validación de reglas de negocio y Excel
avanzado. Pertenece al perfil de propietario de datos, responsable de las definiciones, la
calidad, la vigencia y las reglas de uso de los conjuntos que administra.

\textbf{\textcolor{uxnavy}{Necesidades y objetivos.}} Asegurar que toda la institución use la
definición oficial y la versión más reciente de los datos de crédito. Documentar los metadatos
en un solo lugar para dejar de responder las mismas dudas por correo. Poder aprobar o rechazar
quién consulta el nivel más granular y confidencial de su información.

\textbf{\textcolor{uxnavy}{Pain points relevantes para esta tarea.}} Pierde horas cada semana
explicando a analistas de otras áreas qué significa una columna. Otras áreas siguen trabajando
con extracciones descargadas meses atrás y no se enteran cuando la clasificación oficial
cambia. No existe un repositorio único donde él pueda declarar cuál es la fuente certificada
por su subgerencia, y esa ausencia lo convierte en un cuello de botella permanente.

\elemento{Descripción de la tarea}{qué quiere conseguir y pasos para lograrlo}

El comité regulatorio aprobó una modificación en la clasificación de la cartera de crédito: a
partir del corte de agosto, el saldo vencido se reporta separado del saldo total. Roberto tiene
que publicar la nueva versión del catálogo y retirar la anterior.

Su preocupación no es publicar, es el efecto de publicar. Sabe que al menos una decena de
analistas tiene guardadas consultas y tableros construidos sobre la definición anterior, y que
si la cambia sin avisar, esos análisis van a producir cifras que ya no significan lo mismo. En
la práctica esa situación se detecta meses después, en una reunión donde dos áreas presentan
números distintos. Su meta es que el cambio llegue antes que sus consecuencias.

\figuraescena{imagenes/escena_roberto.png}
  {Contexto del escenario 4: el gobierno del dato es un trabajo de revisión minuciosa cuyo
   efecto se mide en el trabajo de otros. Ilustración generada con inteligencia
   artificial.}{fig:escena-roberto}

\textbf{\textcolor{uxnavy}{Pasos para lograrlo.}}

\begin{uxpreg}
\item Redactar la nueva definición y las reglas de uso que la acompañan.
\item Determinar desde qué fecha de corte aplica y qué ocurre con los periodos anteriores.
\item Identificar qué otras fuentes o indicadores dependen de esta definición.
\item Publicar la nueva versión dejando registro de qué cambió y por qué.
\item Marcar la versión anterior como retirada sin eliminarla.
\item Avisar a quienes usan la fuente, no a toda la institución.
\item Dejar disponible la comparación entre ambas versiones.
\item Poder demostrar después quién fue avisado y cuándo.
\end{uxpreg}

\elemento{Identificación de las interacciones clave}{puntos de contacto entre Roberto y el portal}

Roberto entra al \textbf{módulo de gobierno de metadatos} sobre la ficha de la cartera de
crédito, que él administra. Edita la definición del campo, redacta la regla de uso y añade una
\textbf{nota de conocimiento tribal} explicando el motivo del cambio en lenguaje de negocio,
que es exactamente la explicación que hoy repite por teléfono varias veces al mes.

Declara la \textbf{fecha de vigencia} del cambio. El portal le muestra, antes de publicar, el
\textbf{análisis de impacto}: qué otras fichas del catálogo referencian esta definición, cuántas
consultas guardadas la utilizan y qué tableros la consumen. Es información que hoy no existe en
ningún lado y que él reconstruye preguntando.

Publica la nueva versión y el portal \textbf{marca la anterior como retirada} sin borrarla: queda
consultable, con su periodo de vigencia y un enlace a la vigente. Cualquiera que abra la versión
retirada ve de inmediato que lo es.

La \textbf{notificación dirigida} sale hacia las personas que tienen esa fuente guardada en su
espacio, no hacia una lista de distribución general. El aviso dice qué cambió, desde cuándo
aplica y qué hacer con los análisis producidos antes del cambio. El portal deja
\textbf{registro de la publicación y del aviso}, con fecha y destinatarios.

\elemento{Descripción del resultado deseado}{resultado de la tarea y sensación al completarla}

La nueva clasificación queda publicada, la anterior queda accesible y marcada, y las personas
que dependían de esa definición se enteraron antes de volver a usarla. Roberto ya no tiene que
explicar el cambio de forma individual: la explicación vive junto al dato.

La sensación buscada es \textbf{tranquilidad institucional}: dejar de ser el cuello de botella
y, al mismo tiempo, dejar de ser el responsable implícito de los errores que otros cometen con
sus datos. Es la contraparte exacta del escenario 1: lo que Roberto publica aquí es lo que
Laura lee cuando valida una cifra, y lo que evita que ella tenga que escribirle.

\elemento{Excepciones}{situaciones poco frecuentes con consecuencias de diseño}

\begin{uxlist}
\lead{El cambio afecta a un tablero directivo}{el análisis de impacto lo detecta y exige una
  confirmación adicional antes de publicar, porque una cifra que cambia de significado en un
  tablero de decisión tiene consecuencias distintas a una que cambia en una consulta operativa.}
\lead{Alguien pide acceso al detalle confidencial durante la ventana de cambio}{la solicitud
  queda vinculada a la versión vigente en ese momento, de modo que la autorización que Roberto
  otorgue diga con precisión sobre qué definición se concedió.}
\lead{Hay que revertir la publicación}{la versión anterior sigue existiendo y puede volver a
  marcarse como vigente, con un nuevo registro que documente la reversión. Nada se pierde y
  todo queda explicado.}
\end{uxlist}

% ------------------------------------------------------ Journey individual 2
\clearpage
\section{Journey map individual · Arturo Castañeda}

\textit{Elaborado por Alexandro Mayoral Terán sobre el escenario 3.}

Este mapa despliega el objetivo del escenario 3: \textbf{entender una señal de riesgo y poder
responder por ella ante la Junta de Gobierno}. Es el recorrido más corto de los cuatro —cinco
etapas en quince minutos— y por eso mismo el más exigente: no hay margen para un paso de más.

\textbf{Qué revela el mapa que el escenario no mostraba.} Que el recorrido de un perfil
directivo no es una versión simplificada del recorrido operativo, sino uno distinto. Laura
busca, compara y valida; Arturo no busca nada: la información ya está en su vista y su trabajo
consiste en interpretarla y en decidir si puede responder por ella. Un diseño que tratara al
directivo como a un operativo con menos opciones fallaría en la etapa 3, cuando detecta el
movimiento, porque ahí no necesita más datos sino una explicación.

\textbf{Momento de la verdad.} La etapa 4, cuando toca la gráfica y aparece la tarjeta de
contexto escrita por el propietario del dato. Esa tarjeta sustituye las tres llamadas que hoy
tendría que hacer y es lo único que separa la seguridad de la duda quince minutos antes de una
Junta. Es también el punto donde este recorrido se apoya en el trabajo descrito en el
escenario 4: sin alguien que escriba el contexto, la tarjeta está vacía.

\textbf{Consecuencia para el diseño.} La vista directiva tiene que estar preparada para
lectura y no para exploración. Tres de sus cinco oportunidades —carga instantánea, pocas cifras
con su variación, contexto escrito junto al dato— apuntan a reducir decisiones, no a añadir
capacidades.

\figurapanoramica{figuras/journey_arturo.pdf}
  {Journey map individual de Arturo Castañeda. Objetivo: entender una señal de riesgo y poder
   responder por ella ante la Junta de Gobierno. Elaboración propia.}{fig:journey-arturo}

% ============================================================================
\clearpage
\aportacion[Escenarios]{Arthur Jafed Zizumbo Velasco}{Matrícula A01796363 · Escenarios 5 y 6 · Personas: Mariana Ovalle Ríos y Ximena Solís Barrera}

% ---------------------------------------------------------------- Escenario 5
\section{Escenario 5 · Otorgar el acceso mínimo suficiente y poder retirarlo con la misma rapidez}

\escenariohead{imagenes/persona_07_mariana.png}
  {Otorgar el acceso mínimo suficiente y poder retirarlo con la misma rapidez}
  {Mariana Ovalle Ríos · Perfil administración de la plataforma}
  {Rol en el sistema: administrador · Escritorio · Urgencia alta el día del alta, recurrente el resto del trimestre}

\elemento{Definición del usuario}{audiencia objetivo, necesidades y objetivos, pain points}

Mariana Ovalle Ríos, 41 años, administradora de plataformas de información, ingeniera en
sistemas con especialidad en seguridad de la información y dieciséis años administrando
sistemas institucionales y controles de acceso. Nivel técnico avanzado. Pertenece al perfil de
administración de la plataforma, que gestiona el ciclo de vida de las cuentas: altas, cambios
de rol y bajas.

\textbf{\textcolor{uxnavy}{Necesidades y objetivos.}} Dar acceso el primer día a quien se
incorpora y retirarlo el mismo día de la baja, sin depender de que alguien se acuerde. Asignar
el rol más acotado que permita a cada persona hacer su trabajo, y poder demostrar por qué se
otorgó. Anticipar los problemas de continuidad antes de que los usuarios los reporten.

\textbf{\textcolor{uxnavy}{Pain points relevantes para esta tarea.}} Las cuentas huérfanas
—activas después de que la persona cambió de área o dejó la institución— suelen detectarse
hasta que alguien las audita. Recibe solicitudes de acceso por correo, sin justificación, sin
vigencia y sin visto bueno del propietario del dato. Y su trabajo es invisible: cuando la
plataforma funciona nadie lo nota, y cuando falla no tiene con qué mostrar qué ocurrió.

\elemento{Descripción de la tarea}{qué quiere conseguir y pasos para lograrlo}

Se incorpora una analista a la Dirección de Riesgos y necesita acceso desde su primer día. Un
mes después, la misma persona se mueve a un área directiva. Al cierre del trimestre, Mariana
tiene que revisar todas las cuentas activas y desactivar las que ya no corresponden, con
evidencia lista por si auditoría la solicita.

La tarea es recurrente y aparentemente trivial, pero concentra el riesgo de acceso de toda la
plataforma. Su meta no es dar de alta a una persona: es que en cualquier momento del trimestre
pueda demostrar que cada cuenta activa tiene una razón vigente para estarlo, y que cada baja
ocurrió cuando debía ocurrir.

\figuraescena{imagenes/escena_mariana.png}
  {Contexto del escenario 5: la administración de accesos es un trabajo recurrente y poco
   visible cuyo valor solo se nota cuando falla. Ilustración generada con inteligencia
   artificial.}{fig:escena-mariana}

\textbf{\textcolor{uxnavy}{Pasos para lograrlo.}}

\begin{uxpreg}
\item Recibir la solicitud con la justificación y el visto bueno del responsable.
\item Determinar el rol más acotado que permite hacer el trabajo solicitado.
\item Verificar qué fuentes quedan habilitadas con ese rol antes de confirmarlo.
\item Registrar quién autorizó el acceso y con qué vigencia.
\item Ajustar el rol cuando la persona cambia de función, sin recrear la cuenta.
\item Revisar periódicamente las cuentas activas contra la plantilla vigente.
\item Desactivar las cuentas que sobran, conservando su historial.
\item Obtener la evidencia de todo lo anterior en un formato presentable ante auditoría.
\end{uxpreg}

\elemento{Identificación de las interacciones clave}{puntos de contacto entre Mariana y el portal}

Mariana entra al \textbf{panel de administración} y da de alta a la analista con una contraseña
temporal. Al seleccionar el rol, el portal le muestra en la misma pantalla \textbf{qué fuentes y
qué módulos quedan habilitados} con ese rol: la decisión deja de tomarse a ciegas y ella puede
elegir el permiso mínimo suficiente con la consecuencia a la vista.

El alta queda \textbf{vinculada a la autorización} que la originó: quién la solicitó, quién dio
el visto bueno y con qué vigencia. Ese vínculo es la diferencia entre un registro técnico y
evidencia utilizable.

Un mes después, la persona cambia a un área directiva. Mariana \textbf{modifica el rol} sin
recrear la cuenta; el portal reajusta de inmediato los módulos visibles y \textbf{deja
constancia} de quién hizo el cambio y cuándo. El historial de la cuenta muestra la secuencia
completa.

Al cierre del trimestre, Mariana consulta el \textbf{reporte de cuentas sin actividad}, revisa
las que aparecen y \textbf{desactiva seis en una sola sesión}. La desactivación es lógica: la
cuenta deja de operar pero su historial permanece, de modo que una auditoría posterior puede
reconstruir qué hizo esa persona mientras tuvo acceso. Descarga el \textbf{reporte de accesos
del trimestre} sin pedírselo a nadie.

\elemento{Descripción del resultado deseado}{resultado de la tarea y sensación al completarla}

La analista trabajó desde su primer día con el permiso correcto, el cambio de área no dejó
privilegios residuales y el trimestre cerró con seis cuentas menos y con la evidencia
disponible antes de que alguien la pidiera.

La sensación buscada es \textbf{respaldo}: que el sistema deje el rastro por ella. Es una
necesidad distinta a la eficiencia, y el diseño tiene que distinguirlas. Para Mariana, un panel
que permita dar de alta en diez segundos pero no registre quién autorizó qué le ahorra un
minuto hoy y le cuesta un hallazgo de auditoría en seis meses.

\elemento{Excepciones}{situaciones poco frecuentes con consecuencias de diseño}

\begin{uxlist}
\lead{Solicitud urgente sin justificación}{el portal permite otorgar un acceso temporal con
  vigencia obligatoria y caducidad automática, en lugar de forzar la elección entre bloquear el
  trabajo de alguien o crear un permiso permanente indocumentado.}
\lead{El último administrador intenta desactivarse a sí mismo}{el sistema lo impide y explica
  por qué. Una plataforma sin administrador activo no es un caso raro tolerable: es una
  interrupción de servicio.}
\lead{La persona cambia de área pero conserva funciones del área anterior}{el portal admite la
  situación de forma explícita, con una vigencia declarada para el permiso heredado, de modo que
  la excepción sea visible y caduque sola en lugar de convertirse en una cuenta huérfana.}
\end{uxlist}

% ---------------------------------------------------------------- Escenario 6
\clearpage
\section{Escenario 6 · Integrar un servicio de datos y sobrevivir a un cambio de esquema}

\escenariohead{imagenes/persona_08_ximena.png}
  {Integrar un servicio de datos y sobrevivir a un cambio de esquema}
  {Ximena Solís Barrera · Perfil integración de aplicaciones, consumo programático}
  {Rol en el sistema: analista con credenciales de integración · Escritorio y modalidad híbrida · Urgencia media, ciclo de varias semanas}

\elemento{Definición del usuario}{audiencia objetivo, necesidades y objetivos, pain points}

Ximena Solís Barrera, 28 años, ingeniera de aplicaciones internas con cinco años desarrollando
herramientas y servicios de datos. Nivel técnico avanzado: Python, servicios REST, control de
versiones y automatización. Pertenece al perfil de integración de aplicaciones, que construye
herramientas internas alimentadas por los datos del portal.

\textbf{\textcolor{uxnavy}{Necesidades y objetivos.}} Consumir datos mediante servicios con
una definición documentada que no cambie sin aviso ni versión previa. Eliminar los pasos
manuales de descarga y transformación que hoy repite cada semana. Obtener credenciales,
revisar el catálogo y probar una consulta sin abrir un ticket ni esperar días por una
aprobación.

\textbf{\textcolor{uxnavy}{Pain points relevantes para esta tarea.}} Un campo cambia de nombre
o de tipo y su proceso falla en producción sin que nadie lo hubiera anunciado; la llamada la
recibe ella. Encuentra el servicio pero no la definición del campo, su periodicidad ni sus
valores válidos. Y cada fuente nueva implica un trámite distinto, con formato, responsable y
tiempo de respuesta diferentes, lo que empuja a las áreas a resolver por fuera del canal
oficial.

\elemento{Descripción de la tarea}{qué quiere conseguir y pasos para lograrlo}

El área de riesgos le pide una herramienta que combine la exposición de créditos y derivados
por contraparte, actualizada cada mañana antes de las ocho. Hoy ese cruce lo hace una persona a
mano, cada lunes, descargando dos archivos y pegándolos en una hoja de cálculo.

El reto de Ximena no es escribir el proceso: es construirlo sobre algo que no se rompa. Su
experiencia le dice que el riesgo no está en la primera versión sino en el mes cuatro, cuando
alguien modifique un campo sin avisar y el proceso falle de madrugada. Su meta es entregar una
automatización que sobreviva a los cambios del origen, o que al menos falle de forma anunciada
y no sorpresiva.

\figuraescena{imagenes/escena_ximena.png}
  {Contexto del escenario 6: la integración se construye una vez y se sostiene durante meses;
   el riesgo aparece cuando el origen cambia. Ilustración generada con inteligencia
   artificial.}{fig:escena-ximena}

\textbf{\textcolor{uxnavy}{Pasos para lograrlo.}}

\begin{uxpreg}
\item Localizar en el catálogo los dos conjuntos de datos requeridos.
\item Revisar la definición, la periodicidad y los valores válidos de cada campo.
\item Confirmar qué campo permite relacionar ambos conjuntos.
\item Obtener credenciales de consumo programático para su aplicación.
\item Probar la consulta con un rango corto y verificar el formato de salida.
\item Comprobar qué ocurre cuando el origen no está disponible.
\item Programar la ejecución diaria dentro de la ventana requerida.
\item Suscribirse a los avisos de cambio de esquema de las fuentes que consume.
\item Documentar el contrato de datos sobre el que construyó la herramienta.
\end{uxpreg}

\elemento{Identificación de las interacciones clave}{puntos de contacto entre Ximena y el portal}

Ximena entra al \textbf{catálogo}, localiza los dos conjuntos y revisa las fichas: definición
por campo, periodicidad, valores válidos y propietario. Esa lectura, que hoy exigiría dos
correos y una semana, ocurre en diez minutos y es la que determina si el proyecto es viable.

Genera sus \textbf{credenciales de integración desde su propio perfil}, sin abrir un ticket. El
portal las emite con el mismo alcance que su rol: la aplicación que ella construya no podrá
acceder a nada que ella no pueda ver, y ese límite queda declarado en la propia credencial.

Prueba la consulta con un \textbf{rango corto} y confirma el formato de salida y el campo de
relación entre ambos conjuntos. El portal expone el \textbf{contrato de datos} de la consulta:
qué campos devuelve, de qué tipo y con qué versión de esquema. Ximena guarda ese contrato en su
repositorio y escribe una prueba que detecta cualquier desviación.

Programa la \textbf{ejecución diaria} y se \textbf{suscribe a los avisos de cambio} de las dos
fuentes. Semanas después llega el aviso: uno de los campos cambia de nombre en la próxima
versión. El aviso incluye la fecha de vigencia y confirma que la \textbf{versión anterior
seguirá disponible} durante el periodo de transición. El proceso de Ximena sigue corriendo
mientras ella migra sin prisa.

\elemento{Descripción del resultado deseado}{resultado de la tarea y sensación al completarla}

La herramienta entrega el cruce actualizado cada mañana antes de las ocho, el trabajo manual
del lunes desapareció y el cambio de esquema no interrumpió el servicio en ningún momento.

La sensación buscada es \textbf{confianza en el proveedor de datos}. Para este perfil el
producto no compite contra la ausencia de datos, sino contra la costumbre de descargar
archivos y guardarlos localmente, que es exactamente la práctica que fragmenta la información
institucional. Un portal que no ofrezca un contrato estable no elimina esa costumbre: la
justifica.

\elemento{Excepciones}{situaciones poco frecuentes con consecuencias de diseño}

\begin{uxlist}
\lead{El origen no responde durante la ventana de ejecución}{el servicio responde con un estado
  explícito y una indicación de reintento, en lugar de devolver un resultado vacío que la
  herramienta interpretaría como ausencia de operaciones.}
\lead{El cambio de esquema es inevitable y no admite transición}{el aviso lo declara con esa
  condición y con la mayor anticipación posible, y el portal ofrece la lista de consumidores
  afectados al propietario del dato, de modo que la decisión se tome conociendo su costo.}
\lead{La credencial caduca}{el aviso llega antes del vencimiento y no después del fallo. Una
  credencial que expira en silencio produce exactamente el incidente de madrugada que este
  escenario busca evitar.}
\end{uxlist}

% ------------------------------------------------------ Journey individual 3
\clearpage
\section{Journey map individual · Ximena Solís Barrera}

\textit{Elaborado por Arthur Jafed Zizumbo Velasco sobre el escenario 6.}

Este mapa despliega el objetivo del escenario 6: \textbf{automatizar un cruce diario por
interfaz de programación y sobrevivir a un cambio de esquema}. Es el recorrido más largo en el
tiempo: la etapa final ocurre semanas después de que el trabajo se dio por terminado.

\textbf{Qué revela el mapa que el escenario no mostraba.} Que en este perfil el después del uso
no es un epílogo, sino la etapa que decide todo. Las cinco primeras etapas se resuelven en dos
días y ninguna es especialmente difícil; el valor del producto se juega en la sexta, cuando el
origen cambia. Un mapa que hubiera terminado en la puesta en marcha —como termina la mayoría de
los journey maps— habría dejado fuera precisamente lo que determina si el portal se adopta o se
abandona.

\textbf{Momento de la verdad.} La etapa 4, cuando el portal expone el contrato de datos con su
versión de esquema. Ese contrato es lo que permite a Ximena escribir una prueba que detecta el
cambio antes de que rompa producción. Sin él, las etapas 5 y 6 dependen de la suerte.

\textbf{Consecuencia para el diseño.} El competidor de este perfil no es otra herramienta: es
la costumbre de descargar archivos y guardarlos localmente, que es exactamente la práctica que
fragmenta la información institucional y que el portal existe para revertir. Un portal sin
contrato estable ni aviso anticipado no elimina esa costumbre; la justifica. Es el argumento
más fuerte a favor de tratar la comunicación de cambios como una función del producto y no como
una cortesía.

\figurapanoramica{figuras/journey_ximena.pdf}
  {Journey map individual de Ximena Solís Barrera. Objetivo: automatizar un cruce diario por
   interfaz de programación y sobrevivir a un cambio de esquema. Elaboración
   propia.}{fig:journey-ximena}

